% concatenated from ak_complementary_mccarthy_arxiv.tex
\documentclass[11pt]{amsart}

\usepackage[a4paper,margin=3cm]{geometry}
\usepackage{amsmath,amssymb,mathtools}
\usepackage{enumitem}
\usepackage[hidelinks]{hyperref}
\hypersetup{
  pdftitle={A negative exponent range for Audenaert's complementary McCarthy trace inequality},
  pdfauthor={Xing Li and Bin Zhou}
}

\numberwithin{equation}{section}

\newtheorem{theorem}{Theorem}[section]
\newtheorem{proposition}[theorem]{Proposition}
\newtheorem{lemma}[theorem]{Lemma}
\newtheorem{corollary}[theorem]{Corollary}
\newtheorem{remark}[theorem]{Remark}

\newcommand{\Tr}{\operatorname{Tr}}
\newcommand{\diag}{\operatorname{diag}}
\newcommand{\rank}{\operatorname{rank}}
\newcommand{\norm}[1]{\left\lVert #1\right\rVert}
\newcommand{\abs}[1]{\left\lvert #1\right\rvert}
\newcommand{\C}{\mathbb{C}}

\title[A negative exponent range for a complementary McCarthy inequality]
{A negative exponent range for Audenaert's complementary McCarthy trace inequality}

\author{Xing Li}
\address{Center for Control Theory and Guidance Technology, Harbin Institute of
Technology, Harbin 150001, China}
\email{pipipi1919810@gmail.com}

\author{Bin Zhou}
\address{Center for Control Theory and Guidance Technology, Harbin Institute of
Technology, Harbin 150001, China}
\email{binzhou@gmail.com; binzhou@hit.edu.cn}

\thanks{This paper was supported in part by the National Natural Science
Foundation of China for Distinguished Young Scholars under Grant 62125303; by
the National Natural Science Foundation of China -- ``Qisun Ye'' Science
Foundation under Grant U2441243; by the National Natural Science Foundation of
China under Grants 62573165 and 62521005; by the Natural Science Foundation of
Heilongjiang Province under Grant FG2025F001; and by the Fundamental and
Interdisciplinary Disciplines Breakthrough Plan of the Ministry of Education of
China JYB2025XDXM206.}

\subjclass[2020]{Primary 15A45; Secondary 47A64, 47B10}
\keywords{McCarthy inequality, trace inequality, Schatten norm, Kubo--Ando mean, parallel sum, log-majorization}

\begin{document}

\begin{abstract}
Audenaert introduced a class of complementary McCarthy type trace inequalities in his work on completely monotone functions and Bernstein functions; the same problem was later included in the problem list of Audenaert and Kittaneh.  The known results cover the corresponding directions for \(q\le -2\), \(0<q\le 1\), \(1\le q\le 2\), and \(2\le q\le 3\), while the negative range \(-2<q<0\) was left as a conjectural case.  We prove the negative exponent inequality, in fact for every \(q<0\).  After the inversion \(X=A^{-1}\), \(Y=B^{-1}\), the problem reduces to a trace inequality for the parallel sum \(X:Y\).  The proof uses the Kubo--Ando mean chain, an Ando--Hiai type log-majorization for matrix geometric means, and a finite-dimensional Schatten H\"older inequality, including the quasi-norm range.  We also record that the positive exponent side \(q\ge 3\) follows directly from Audenaert's norm-compression inequality for positive semidefinite \(2\times2\) block matrices, and in fact holds for all \(q\ge 2\).  Finally, we determine the equality case on the negative side: equality holds if and only if \(A=B\).
\end{abstract}

\maketitle

\section{Introduction}

Let \(A,B\) be positive matrices.  McCarthy's trace inequality says that
\begin{equation}\label{eq:mccarthy}
    \Tr(A+B)^q\ge \Tr A^q+\Tr B^q,\qquad q\ge 1,
\end{equation}
while the reverse inequality holds for \(0\le q\le 1\); see McCarthy's work on Schatten classes \cite{McCarthy1967}.  In his study of trace inequalities for completely monotone functions and Bernstein functions, Audenaert obtained complementary estimates involving the mixed term
\[
    \Tr(A^{1/2}BA^{1/2})^{q/2}.
\]
These estimates are not merely perturbations of \eqref{eq:mccarthy}.  In the positive exponent range they strengthen McCarthy's inequality, whereas in the negative exponent range they give nontrivial estimates for inverse powers.

The complementary McCarthy type inequalities in question can be written as follows.  For \(q<0\), and for positive definite \(A,B\), the expected direction is
\begin{equation}\label{eq:direct-q}
    \Tr(A+B)^q
    \le
    \Tr A^q+\Tr B^q
    +(2^q-2)\Tr(A^{1/2}BA^{1/2})^{q/2}.
\end{equation}
For positive exponents, the corresponding reverse inequality is
\begin{equation}\label{eq:reverse-q}
    \Tr(A+B)^q
    \ge
    \Tr A^q+\Tr B^q
    +(2^q-2)\Tr(A^{1/2}BA^{1/2})^{q/2}.
\end{equation}
Audenaert's Corollary 3 in \cite{Audenaert2012Trace} proves \eqref{eq:direct-q} for \(q\le -2\) and \(1\le q\le 2\), and proves \eqref{eq:reverse-q} for \(0<q\le 1\) and \(2\le q\le 3\).  The problem list of Audenaert and Kittaneh \cite[Section 3]{AudenaertKittaneh2012Problems} then singled out the range \(-2<q<0\) and the positive range \(q\ge 3\) as conjectural.  The first of these is the main point of the present paper.  The second is best viewed as a consequence of an earlier result: Audenaert's norm-compression inequality for positive semidefinite \(2\times2\) block matrices \cite{Audenaert2006NormCompression} directly implies \eqref{eq:reverse-q} for all \(q\ge 2\).

Our main result is the following.

\begin{theorem}\label{thm:main}
Let \(A,B\) be positive matrices.
\begin{enumerate}[label=\textup{(\roman*)}]
    \item If \(A,B\) are positive definite, then \eqref{eq:direct-q} holds for every \(q<0\).
    \item If \(A,B\) are positive semidefinite, then \eqref{eq:reverse-q} holds for every \(q\ge 2\).
\end{enumerate}
Thus the negative exponent range \(-2<q<0\) in the Audenaert--Kittaneh problem list is settled, while the positive exponent range \(q\ge 3\) is covered, and strengthened to \(q\ge 2\), by a known norm-compression inequality.
\end{theorem}

The positive exponent part of Theorem \ref{thm:main} is a bibliographical observation rather than a new block matrix inequality.  We include the proof in Section \ref{sec:positive}, because the substitution is short and clarifies the natural endpoint \(q=2\).

The negative exponent part is the new contribution.  Put \(q=-r\) and make the inversion
\[
    X=A^{-1},\qquad Y=B^{-1}.
\]
Then
\[
    (A+B)^{-1}=X:Y=(X^{-1}+Y^{-1})^{-1},
\]
where \(X:Y\) denotes the parallel sum.  The main estimate is a trace inequality for \(X:Y\).  We first prove a one-parameter asymmetric form in which the mixed term
\[
    \norm{X^{1/2}Y^{1/2}}_r^r
\]
is replaced by
\[
    \norm{X^{1-t}Y^t}_r^r,\qquad 0<t<1.
\]
The symmetric value \(t=1/2\) gives exactly the form needed for \eqref{eq:direct-q}.  We do not claim that the one-parameter constants are optimal when \(t\ne 1/2\); the sharp case needed here is the symmetric one.

All matrices in the paper are finite-dimensional complex matrices.  The trace is the usual non-normalized trace.  For \(p>0\) and a matrix \(Z\), we write
\[
    \norm{Z}_p=(\Tr |Z|^p)^{1/p}.
\]
When \(0<p<1\), this is the Schatten \(p\)-quasi-norm.

\section{Preliminaries}\label{sec:prelim}

For \(X,Y>0\), the parallel sum is
\[
    X:Y=(X^{-1}+Y^{-1})^{-1}.
\]
We use the Kubo--Ando harmonic and geometric means
\[
    X!Y=2(X:Y),\qquad
    X\#Y=X^{1/2}(X^{-1/2}YX^{-1/2})^{1/2}X^{1/2}.
\]
More generally, for \(0<t<1\), set
\[
    X!_tY=((1-t)X^{-1}+tY^{-1})^{-1},
    \qquad
    X\#_tY=X^{1/2}(X^{-1/2}YX^{-1/2})^tX^{1/2}.
\]
These means satisfy the standard weighted mean chain
\begin{equation}\label{eq:weighted-mean-chain}
    X!_tY\le X\#_tY\le (1-t)X+tY.
\end{equation}
At \(t=1/2\) this reduces to
\begin{equation}\label{eq:mean-chain}
    X!Y\le X\#Y\le \frac{X+Y}{2}.
\end{equation}
This is the standard order relation between the harmonic, geometric, and arithmetic Kubo--Ando means; we use it only for finite-dimensional positive definite matrices.  See Kubo--Ando \cite{KuboAndo1980} and Bhatia \cite{Bhatia2007}.

For a matrix \(Z\), let \(s(Z)=(s_1(Z),\ldots,s_n(Z))\) be its singular values in decreasing order.  For a positive matrix \(W\), let \(\lambda(W)\) be its eigenvalue vector in decreasing order.  We write \(\prec_{\log}\) for log-majorization.

\begin{lemma}[Geometric means and products]\label{lem:ando-hiai}
Let \(X,Y>0\) and \(0<t<1\).  Then
\begin{equation}\label{eq:AH-weighted}
    \lambda(X\#_tY)\prec_{\log}s(X^{1-t}Y^t).
\end{equation}
Consequently, for every \(r>0\),
\begin{equation}\label{eq:geomean-trace-bound-weighted}
    \Tr(X\#_tY)^r\le \norm{X^{1-t}Y^t}_r^r.
\end{equation}
In particular, when \(t=1/2\),
\begin{equation}\label{eq:geomean-trace-bound}
    \Tr(X\#Y)^r\le \norm{X^{1/2}Y^{1/2}}_r^r.
\end{equation}
\end{lemma}

\begin{proof}
The log-majorization \eqref{eq:AH-weighted} is a standard form of the Ando--Hiai log-majorization theory for matrix geometric means.  Our convention for the weighted geometric mean is
\[
    X\#_tY=X^{1/2}(X^{-1/2}YX^{-1/2})^tX^{1/2},
\]
and the right-hand side of \eqref{eq:AH-weighted} denotes the singular values of the product \(X^{1-t}Y^t\).  This form follows from the log-majorization results of Ando--Hiai \cite{AndoHiai1994}; see also Bhatia \cite[Chapter IX]{Bhatia1997} for the discussion of geometric means and log-majorization.  We only use the finite-dimensional positive definite case.

If \(x\prec_{\log}y\) and the entries are non-negative and arranged in decreasing order, then for every \(r>0\), the vector \((x_1^r,\ldots,x_n^r)\) is weakly majorized by \((y_1^r,\ldots,y_n^r)\).  Hence
\[
    \sum_j \lambda_j(X\#_tY)^r
    \le
    \sum_j s_j(X^{1-t}Y^t)^r,
\]
which is \eqref{eq:geomean-trace-bound-weighted}.  Taking \(t=1/2\) gives \eqref{eq:geomean-trace-bound}.
\end{proof}

\begin{lemma}[Trace monotonicity and equality]\label{lem:trace-monotone}
If \(0\le C\le D\), then for every \(r>0\),
\[
    \Tr C^r\le \Tr D^r.
\]
Moreover, if equality holds, then \(C=D\).
\end{lemma}

\begin{proof}
By the min-max principle, the eigenvalues in decreasing order satisfy
\[
    \lambda_j(C)\le \lambda_j(D),\qquad j=1,\ldots,n.
\]
Taking the positive power \(r\) and summing gives \(\Tr C^r\le \Tr D^r\).

If \(\Tr C^r=\Tr D^r\), then every scalar inequality
\[
    \lambda_j(C)^r\le \lambda_j(D)^r
\]
must be an equality.  Thus \(\Tr C=\Tr D\).  Since \(D-C\ge 0\), we get
\[
    \Tr(D-C)=0,
\]
and hence \(D-C=0\).
\end{proof}

\begin{lemma}[Finite-dimensional Schatten H\"older inequality]\label{lem:holder}
Let \(p,q,r>0\) and \(1/r=1/p+1/q\).  For all compatible matrices \(U,V\),
\[
    \norm{UV}_r\le \norm{U}_p\norm{V}_q.
\]
In particular, for \(X,Y>0\), \(r>0\), and \(0<t<1\),
\begin{equation}\label{eq:holder-weighted-special}
    \norm{X^{1-t}Y^t}_r^r
    \le
    (\Tr X^r)^{1-t}(\Tr Y^r)^t.
\end{equation}
At \(t=1/2\),
\begin{equation}\label{eq:holder-special}
    \norm{X^{1/2}Y^{1/2}}_r^r
    \le
    (\Tr X^r)^{1/2}(\Tr Y^r)^{1/2}
    \le
    \frac{\Tr X^r+\Tr Y^r}{2}.
\end{equation}
\end{lemma}

\begin{proof}
We first prove the general form.  In finite dimensions, Horn's singular value product inequality gives the weak log-majorization
\[
    s(UV)\prec_{w\log} s(U)s(V),
\]
where the right-hand side is the componentwise product of the singular value vectors, both arranged in decreasing order; see Bhatia \cite[Chapter III]{Bhatia1997}.  Raising entries to the power \(r\), summing, and using the summation property of weak majorization gives
\[
    \norm{UV}_r^r
    =\sum_j s_j(UV)^r
    \le
    \sum_j s_j(U)^r s_j(V)^r.
\]
Since \(1/r=1/p+1/q\), we have \(p/r>1\), \(q/r>1\), and
\[
    \frac{1}{p/r}+\frac{1}{q/r}=1.
\]
The scalar H\"older inequality yields
\[
    \sum_j s_j(U)^r s_j(V)^r
    \le
    \left(\sum_j s_j(U)^p\right)^{r/p}
    \left(\sum_j s_j(V)^q\right)^{r/q}
    =
    \norm{U}_p^r\norm{V}_q^r.
\]
Taking \(r\)-th roots proves \(\norm{UV}_r\le \norm{U}_p\norm{V}_q\).  No triangle inequality for \(\norm{\cdot}_p\) is used, so \(p,q,r\) may be smaller than \(1\).

Now take \(U=X^{1-t}\), \(V=Y^t\), \(p=r/(1-t)\), and \(q=r/t\).  Then
\[
    \norm{X^{1-t}Y^t}_r
    \le
    \norm{X^{1-t}}_{r/(1-t)}
    \norm{Y^t}_{r/t}.
\]
Since \(X,Y\) are positive,
\[
    \norm{X^{1-t}}_{r/(1-t)}=(\Tr X^r)^{(1-t)/r},
    \qquad
    \norm{Y^t}_{r/t}=(\Tr Y^r)^{t/r}.
\]
Raising to the power \(r\) gives \eqref{eq:holder-weighted-special}.  The first inequality in \eqref{eq:holder-special} is the case \(t=1/2\), and the second one is the scalar arithmetic-geometric mean inequality.
\end{proof}

We also need the following norm-compression inequality of Audenaert.

\begin{lemma}[Audenaert's norm-compression inequality]\label{lem:compression}
Let
\[
    M=\begin{pmatrix}P&Q\\ Q^*&R\end{pmatrix}\ge 0,
\]
where \(P\) and \(R\) are square diagonal blocks.  If \(p\ge 2\), then
\begin{equation}\label{eq:compression}
    \norm{M}_p^p
    \ge
    \norm{P}_p^p+\norm{R}_p^p+(2^p-2)\norm{Q}_p^p.
\end{equation}
For \(1\le p\le 2\), the inequality is reversed.
\end{lemma}

\begin{proof}
This is Audenaert's norm-compression inequality for positive semidefinite \(2\times2\) block matrices \cite[Theorem 1 and Corollary 1]{Audenaert2006NormCompression}.  The hypotheses in the cited result are precisely that the block matrix is positive semidefinite and that the diagonal blocks are square.  Theorem 1 of \cite{Audenaert2006NormCompression} gives the direction \(1\le p\le 2\), while Corollary 1 gives the reverse direction for \(p\ge 2\) by duality.
\end{proof}

\section{A trace inequality for parallel sums}\label{sec:parallel}

The negative exponent inequality will follow from a trace estimate for the parallel sum.  We first prove a one-parameter asymmetric form.  For \(0<t<1\), set
\[
    m_t=\max\{t,1-t\}.
\]

\begin{theorem}[A one-parameter parallel sum inequality]\label{thm:parallel-asym}
Let \(X,Y>0\), \(r>0\), and \(0<t<1\).  Then
\begin{equation}\label{eq:parallel-asym}
    \Tr(X:Y)^r+
    \bigl(m_t^{-1}-m_t^r\bigr)\norm{X^{1-t}Y^t}_r^r
    \le
    \Tr X^r+\Tr Y^r .
\end{equation}
\end{theorem}

\begin{proof}
Put
\[
    H=X:Y,\qquad
    G_t=\norm{X^{1-t}Y^t}_r^r,\qquad
    S=\Tr X^r+\Tr Y^r.
\]
Since
\[
    (1-t)X^{-1}+tY^{-1}
    \le
    m_t(X^{-1}+Y^{-1}),
\]
inverting gives
\[
    X:Y\le m_t(X!_tY).
\]
By \eqref{eq:weighted-mean-chain},
\[
    X:Y\le m_t(X\#_tY).
\]
Using Lemmas \ref{lem:trace-monotone} and \ref{lem:ando-hiai}, we get
\begin{equation}\label{eq:asym-H-estimate}
    \Tr H^r
    \le
    m_t^r\Tr(X\#_tY)^r
    \le
    m_t^r G_t.
\end{equation}
On the other hand, Lemma \ref{lem:holder}, in the form \eqref{eq:holder-weighted-special}, gives
\begin{equation}\label{eq:asym-holder}
    G_t
    =
    \norm{X^{1-t}Y^t}_r^r
    \le
    (\Tr X^r)^{1-t}(\Tr Y^r)^t.
\end{equation}
The scalar weighted arithmetic-geometric mean inequality yields
\[
    (\Tr X^r)^{1-t}(\Tr Y^r)^t
    \le
    (1-t)\Tr X^r+t\Tr Y^r
    \le
    m_tS.
\]
Hence
\begin{equation}\label{eq:asym-G-estimate}
    m_t^{-1}G_t\le S.
\end{equation}
Combining \eqref{eq:asym-H-estimate} and \eqref{eq:asym-G-estimate}, we obtain
\[
    \Tr H^r+
    \bigl(m_t^{-1}-m_t^r\bigr)G_t
    \le
    m_t^rG_t+
    \bigl(m_t^{-1}-m_t^r\bigr)G_t
    =
    m_t^{-1}G_t
    \le S.
\]
This is \eqref{eq:parallel-asym}.
\end{proof}

\begin{remark}
The one-parameter form \eqref{eq:parallel-asym} is used as a natural proof framework for the symmetric case.  We do not claim that the constant is optimal for every \(t\ne 1/2\).
\end{remark}

\begin{theorem}[The symmetric parallel sum inequality]\label{thm:parallel}
Let \(X,Y>0\) and \(r>0\).  Then
\begin{equation}\label{eq:parallel-theorem}
    \Tr(X:Y)^r+
    \left(2-2^{-r}\right)\norm{X^{1/2}Y^{1/2}}_r^r
    \le
    \Tr X^r+\Tr Y^r.
\end{equation}
The coefficient \(2-2^{-r}\) is optimal, in the sense that it cannot be increased uniformly over all positive definite \(X,Y\).
\end{theorem}

\begin{proof}
Take \(t=1/2\) in Theorem \ref{thm:parallel-asym}.  Then \(m_t=1/2\) and
\[
    m_t^{-1}-m_t^r=2-2^{-r},
\]
which gives \eqref{eq:parallel-theorem}.

If \(X=Y\), then \(X:Y=X/2\) and
\[
    \norm{X^{1/2}Y^{1/2}}_r^r=\Tr X^r.
\]
Thus \eqref{eq:parallel-theorem} is an equality at \(X=Y\).  Increasing the coefficient would therefore make the inequality fail at this point.
\end{proof}

\section{The negative exponent range}\label{sec:negative}

We now derive the negative exponent part of Theorem \ref{thm:main}.

\begin{theorem}\label{thm:negative}
Let \(A,B>0\) and \(q<0\).  Then
\[
    \Tr(A+B)^q
    \le
    \Tr A^q+\Tr B^q+
    (2^q-2)\Tr(A^{1/2}BA^{1/2})^{q/2}.
\]
\end{theorem}

\begin{proof}
Write \(q=-r\), where \(r>0\), and set
\[
    X=A^{-1},\qquad Y=B^{-1}.
\]
Then
\begin{equation}\label{eq:AB-parallel}
    (A+B)^{-1}=(X^{-1}+Y^{-1})^{-1}=X:Y.
\end{equation}
Moreover,
\[
    (A^{1/2}BA^{1/2})^{-1}
    =
    A^{-1/2}B^{-1}A^{-1/2}
    =
    X^{1/2}YX^{1/2}.
\]
The positive matrices \(X^{1/2}YX^{1/2}\) and \(Y^{1/2}XY^{1/2}\) have the same nonzero eigenvalues, and
\[
    Y^{1/2}XY^{1/2}
    =
    (X^{1/2}Y^{1/2})^*(X^{1/2}Y^{1/2}).
\]
Therefore
\begin{equation}\label{eq:mixed-transform}
    \Tr(A^{1/2}BA^{1/2})^{-r/2}
    =
    \norm{X^{1/2}Y^{1/2}}_r^r.
\end{equation}
Applying Theorem \ref{thm:parallel} to \(X,Y\), and then using \eqref{eq:AB-parallel} and \eqref{eq:mixed-transform}, gives
\[
    \Tr(A+B)^{-r}
    +
    \left(2-2^{-r}\right)
    \Tr(A^{1/2}BA^{1/2})^{-r/2}
    \le
    \Tr A^{-r}+\Tr B^{-r}.
\]
Since \(q=-r\), this is the desired inequality.
\end{proof}

\section{The positive exponent range}\label{sec:positive}

The positive exponent range is a direct application of Lemma \ref{lem:compression}.  This section is included to make explicit that the range \(q\ge 3\) in the Audenaert--Kittaneh problem list is already covered by the norm-compression inequality, with endpoint \(q=2\).

\begin{proposition}\label{prop:positive}
Let \(A,B\ge 0\) and \(q\ge 2\).  Then
\[
    \Tr(A+B)^q
    \ge
    \Tr A^q+\Tr B^q+
    (2^q-2)\Tr(A^{1/2}BA^{1/2})^{q/2}.
\]
\end{proposition}

\begin{proof}
Let
\[
    V=\begin{pmatrix}A^{1/2}\\ B^{1/2}\end{pmatrix},
    \qquad
    M=VV^*
    =
    \begin{pmatrix}
        A&A^{1/2}B^{1/2}\\
        B^{1/2}A^{1/2}&B
    \end{pmatrix}.
\]
Then \(M=VV^*\ge 0\), and \(M\) is a \(2\times2\) positive semidefinite block matrix whose diagonal blocks \(A,B\) are square.  Hence \(M\) satisfies the hypotheses of Lemma \ref{lem:compression}.  The matrices \(VV^*\) and \(V^*V\) have the same nonzero eigenvalues, while
\[
    V^*V=A+B.
\]
Thus
\begin{equation}\label{eq:V-spectrum}
    \norm{M}_q^q=\Tr M^q=\Tr(A+B)^q.
\end{equation}
Applying Lemma \ref{lem:compression} to \(M\), with
\[
    P=A,\qquad R=B,\qquad Q=A^{1/2}B^{1/2},
\]
we obtain
\begin{equation}\label{eq:positive-step}
    \Tr(A+B)^q
    \ge
    \Tr A^q+\Tr B^q+
    (2^q-2)\norm{A^{1/2}B^{1/2}}_q^q.
\end{equation}
Finally,
\[
    \norm{A^{1/2}B^{1/2}}_q^q
    =
    \Tr(B^{1/2}AB^{1/2})^{q/2}.
\]
The positive semidefinite matrices \(B^{1/2}AB^{1/2}\) and \(A^{1/2}BA^{1/2}\) have the same nonzero eigenvalues.  Hence
\[
    \Tr(B^{1/2}AB^{1/2})^{q/2}
    =
    \Tr(A^{1/2}BA^{1/2})^{q/2}.
\]
Substitution into \eqref{eq:positive-step} proves the proposition.
\end{proof}

\section{Equality and sharpness}\label{sec:sharpness}

We first determine the equality case in the negative exponent range.

\begin{proposition}[Equality for negative exponents]\label{prop:equality-negative}
Let \(A,B>0\) and \(q<0\).  Then
\[
    \Tr(A+B)^q
    =
    \Tr A^q+\Tr B^q+
    (2^q-2)\Tr(A^{1/2}BA^{1/2})^{q/2}
\]
if and only if \(A=B\).
\end{proposition}

\begin{proof}
If \(A=B\), direct substitution gives equality.

Conversely, put \(q=-r\), \(X=A^{-1}\), and \(Y=B^{-1}\).  If equality holds in the displayed formula, then equality must hold in the estimates used to prove Theorem \ref{thm:parallel-asym} at \(t=1/2\).  More explicitly, with
\[
    H=X:Y,\qquad
    G=\norm{X^{1/2}Y^{1/2}}_r^r,\qquad
    S=\Tr X^r+\Tr Y^r,
\]
the proof used
\[
    \Tr H^r\le 2^{-r}\Tr(X\#Y)^r\le 2^{-r}G,
    \qquad
    2G\le S.
\]
Equality in the final result forces equality throughout these two chains.  In particular, from
\[
    H=X:Y\le \frac12(X\#Y)
\]
and Lemma \ref{lem:trace-monotone}, we get
\[
    X:Y=\frac12(X\#Y).
\]
Let \(Z=X^{-1/2}YX^{-1/2}\).  By the Kubo--Ando representations, the last equality is equivalent to
\[
    I:Z=\frac12 Z^{1/2}.
\]
On the spectrum of \(Z\), this says
\[
    \frac{z}{1+z}=\frac12 z^{1/2},\qquad z>0.
\]
Equivalently, \((z^{1/2}-1)^2=0\).  Thus \(Z=I\), so \(X=Y\), and hence \(A=B\).
\end{proof}

The coefficient \(2^q-2\) is forced by the symmetric normalization.  Indeed, if \(A=B>0\), then
\[
    \Tr(A+B)^q=2^q\Tr A^q,
    \qquad
    \Tr(A^{1/2}BA^{1/2})^{q/2}=\Tr A^q.
\]
Therefore
\[
    \Tr A^q+\Tr B^q+
    (2^q-2)\Tr(A^{1/2}BA^{1/2})^{q/2}
    =
    2^q\Tr A^q.
\]
Thus both directions in Theorem \ref{thm:main} are equalities when \(A=B\).  On the negative side, since \(2^q-2<0\), the coefficient cannot be decreased.  On the positive side, the coefficient cannot be increased.

For completeness, we also record the scalar normalization.  For positive numbers \(a,b\), let
\[
    \Phi_q(a,b)
    =
    (a+b)^q-a^q-b^q-(2^q-2)(ab)^{q/2}.
\]
By homogeneity, one may take \(a=e^t\), \(b=e^{-t}\).  Then
\begin{equation}\label{eq:scalar-normalization}
    \Phi_q(e^t,e^{-t})
    =
    2^q(\cosh^q t-1)-2(\cosh(qt)-1).
\end{equation}
Near \(t=0\),
\begin{equation}\label{eq:scalar-expansion}
    \Phi_q(e^t,e^{-t})
    =
    q(2^{q-1}-q)t^2+O(t^4).
\end{equation}
This local expansion is consistent with the signs in Theorem \ref{thm:main}: the quadratic coefficient is negative for \(q<0\) and non-negative for \(q\ge 2\), with equality at \(q=2\).  The proof above is not a reduction to the scalar case; the scalar formula only explains the coefficient and the change of direction.

\section{Conclusion}

Theorem \ref{thm:negative} and Proposition \ref{prop:positive} together prove Theorem \ref{thm:main}.  The positive exponent range is a direct consequence of Audenaert's norm-compression inequality for positive semidefinite block matrices.  The new part is the negative exponent range \(q<0\), in particular the interval \(-2<q<0\) left open in the Audenaert--Kittaneh problem list.  The proof is based on the symmetric case of the parallel sum trace inequality \eqref{eq:parallel-theorem}, and the equality analysis shows that equality in the negative exponent range occurs exactly at \(A=B\).

\begin{thebibliography}{10}

\bibitem{AndoHiai1994}
T.~Ando and F.~Hiai.
\newblock Log majorization and complementary Golden--Thompson type inequalities.
\newblock \emph{Linear Algebra and its Applications}, 197/198:113--131, 1994.

\bibitem{Audenaert2006NormCompression}
K.~M.~R. Audenaert.
\newblock A norm compression inequality for block partitioned positive semidefinite matrices.
\newblock \emph{Linear Algebra and its Applications}, 413(1):155--176, 2006.

\bibitem{Audenaert2012Trace}
K.~M.~R. Audenaert.
\newblock Trace inequalities for completely monotone functions and Bernstein functions.
\newblock \emph{Linear Algebra and its Applications}, 437(2):601--611, 2012.

\bibitem{AudenaertKittaneh2012Problems}
K.~M.~R. Audenaert and F.~Kittaneh.
\newblock Problems and conjectures in matrix and operator inequalities, 2012.
\newblock arXiv:1201.5232.

\bibitem{Bhatia1997}
R.~Bhatia.
\newblock \emph{Matrix Analysis}, volume 169 of \emph{Graduate Texts in Mathematics}.
\newblock Springer, New York, 1997.

\bibitem{Bhatia2007}
R.~Bhatia.
\newblock \emph{Positive Definite Matrices}.
\newblock Princeton Series in Applied Mathematics. Princeton University Press, Princeton, 2007.

\bibitem{KuboAndo1980}
F.~Kubo and T.~Ando.
\newblock Means of positive linear operators.
\newblock \emph{Mathematische Annalen}, 246:205--224, 1980.

\bibitem{McCarthy1967}
C.~A. McCarthy.
\newblock $c_p$.
\newblock \emph{Israel Journal of Mathematics}, 5:249--271, 1967.

\end{thebibliography}

\end{document}
